\documentclass[11pt]{article}
\usepackage[T1]{fontenc}
\usepackage{textcomp}
\usepackage[margin=0.75in]{geometry}
\usepackage{amsmath,amssymb,amsthm}
\usepackage{array,booktabs}
\usepackage{hyperref}
\setlength{\parindent}{0pt}
\newtheorem{theorem}{Theorem}
\theoremstyle{definition}
\newtheorem{definition}{Definition}
\title{Flow Proof Artifact: prop-08-right-triangle-perpendicular-similar}
\author{Generated by \texttt{flow doc proof}}
\date{Flow Proof Book}
\begin{document}
\maketitle
If in a right-angled triangle a perpendicular is drawn from the right angle to the base, the triangles on each side are similar to the whole and to one another.\\[0.5em]
\textbf{Source.} euclid — Elements, Book VI, Proposition 8\\[0.5em]
\bigskip
\noindent{\large \textbf{Derived fact 1.} \textit{If in a right-angled triangle a perpendicular is drawn from the right angle to the base, the triangles on each side are similar to the whole and to one another} \hfill the Euclidean plane \cdot Euclid Book VI \cdot Proposition 8: in a right triangle the altitude to the hypotenuse gives similar subtriangles}\par
\smallskip\noindent\textit{Coordinate: the Euclidean plane \cdot Euclid Book VI \cdot Proposition 8: in a right triangle the altitude to the hypotenuse gives similar subtriangles}\par
\smallskip\noindent\textit{Source: euclid — Elements, Book VI, Proposition 8}\par
\smallskip\noindent\textit{Built on: proposition 47: square on hypotenuse equals sum of squares on legs, for Book I of the Elements in on the Euclidean plane, proposition 4: equiangular triangles have proportional corresponding sides, for Euclid Book VI on the Euclidean plane}\par
\medskip
\noindent\textbf{Goal.} If in a right-angled triangle a perpendicular is drawn from the right angle to the base, the triangles on each side are similar to the whole and to one another\par
\noindent$\displaystyle \text{right triangle altitude gives similarity}$\par
\medskip
\noindent
\begin{tabular}{@{} >{\raggedright\arraybackslash}p{0.06\textwidth} >{\raggedright\arraybackslash}p{0.47\textwidth} >{\raggedleft\arraybackslash}p{0.41\textwidth} @{}}
\textbf{\#} & \textbf{Proof} & \textbf{Mathematics} \\
\midrule
\textbf{1.} & We prove that proposition 8: in a right triangle the altitude to the hypotenuse gives similar subtriangles for Euclid Book VI the Euclidean plane. & \\[0.45em]
\textbf{2.} & Let ABC = right\_triangle. & \\[0.45em]
\textbf{3.} & Let AD = altitude\_to\_hypotenuse. & \\[0.45em]
\textbf{4.} & Let ABD = subtriangle. & \\[0.45em]
\textbf{5.} & From step 2, step 3, and step 4, we can deduce that triangle ABD \textasciitilde{} triangle ABC. & $\displaystyle \text{triangle ABD} ~ \triangle ABC$ \\[0.45em]
\textbf{6.} & We invoke the derived fact governing Book I of the Elements in the Euclidean plane: proposition 47: square on hypotenuse equals sum of squares on legs, for Book I of the Elements in on the Euclidean plane. & \\[0.45em]
\textbf{7.} & We invoke the derived fact governing Euclid Book VI the Euclidean plane: proposition 4: equiangular triangles have proportional corresponding sides, for Euclid Book VI on the Euclidean plane. & \\[0.45em]
\textbf{8.} & From step 6 and step 7, this implies right triangle altitude gives similarity. Hence proven. & $\displaystyle \text{right triangle altitude gives similarity}$ \\[0.45em]
\bottomrule
\end{tabular}
\label{thm:Geometry-Euclid Book VI-proposition_8_in_a_right_triangle_the_altitude_to_the_hypotenuse_gives_similar_subtriangles}
\par\medskip\hrule
\end{document}
